\section{Five Transitions}

A candidate is anything that might be taken to count: a statement, a proof attempt, a model, a policy, a patch, a decision. Whether a candidate counts is rarely settled by a single event. It is usually moved through a sequence of operations, and the operations are easy to confuse because in ordinary speech one word often covers several of them. A claim is said to have been made, published, accepted, approved and adopted as if these were synonyms, when each names a different change in the state of some shared arrangement. This section separates five such changes. They are called transitions because each takes a system from one state to another, and because what is wrong with a given case is usually located in a transition that was assumed to have occurred, or assumed to follow from another, and did not.

Two notions must be kept apart throughout. A transition \emph{occurs} when it is among the things a system actually does, and this is a fact about the history of the system. A transition is \emph{licensed} when the rules of the system permit it in the state in which it is attempted, and this is a fact about the rules. Neither is the same as a transition being \emph{warranted}, which is the subject of the next section and concerns whether it ought to count for a stated consequence on the evidence available. A system can do what its rules forbid, its rules can permit what no evidence supports, and its rules can forbid what the evidence would amply support. Much confusion in discussions of automated systems and of institutions alike comes from reading one of these off another.

\emph{Generation} is the production of a candidate. A language model emitting a proof attempt is a generation, and so is a person drafting a clause or a search procedure proposing a configuration. Generation adds a candidate to someone's workspace and does not, by itself, add anything to the shared record. Its preservation condition is therefore modest: generating must not silently alter the shared record, and the candidate must remain distinguishable from anything the record contains. Its answerability condition is likewise modest and concerns provenance, since whoever or whatever produced the candidate should be identifiable if the candidate later matters.

\emph{Inscription} is the entry of a candidate into a shared record. A document filed in a land registry, a commit pushed to a shared repository, and an entry added to a ledger are inscriptions. What inscription changes is who can see and cite the candidate, and it changes nothing about whether the candidate is correct. Its preservation condition is that the record keep the inscribed item distinguishable from items that have undergone any further transition, so that the presence of an item in the record is not read as evidence that it has been checked. Its answerability condition falls on whoever inscribes, who must be able to answer for the accuracy of what was entered and the provenance attached to it, though not for its truth.

\emph{Evaluation} is the application of an evaluation procedure to a candidate on the basis of evidence, producing an assessed outcome that is itself recorded. A proof checker run on a derivation, a referee report, a test suite executed against a patch and an audit of a ledger entry are evaluations. The preservation condition is the most demanding of the five, because an outcome without a record of what was evaluated, against what evidence and by what procedure cannot later be reassessed, and an assessment that cannot be reassessed has been turned into an assertion. The answerability condition falls on whoever applies and maintains the procedure, who must be able to say what the procedure does and does not establish.

\emph{Authorization} is the permission, by a body with jurisdiction, to rely on a candidate for a stated use. A regulator approving a drug for a specified indication, a court admitting evidence for a specified purpose, and a committee approving a design for production are authorizations. Authorization differs from evaluation in that it is an act of an authority and not an assessment of evidence, and it can be wrong while the evaluation behind it was sound, or right while the evaluation behind it was poor. Its preservation condition is that the use and scope for which permission was given remain attached to the permission, since an authorization stripped of its scope reads as a general warrant. Its answerability condition falls on the authorizing body.

\emph{Commitment} is a change in a system such that reversing it would require further transitions and others have begun to depend on it. A release tag that downstream projects build against, a ratified statute, a signed contract, and a published result on which later work proceeds are commitments. Commitment changes the state in which later transitions are assessed, which is why it receives a separate treatment of its own \cite{admiss}. Its preservation condition is that the structure of dependence be recoverable, meaning that it remain possible to say who relies on what. Its answerability condition falls on those who made the commitment and on those who have relied on it in ways that could not have been withdrawn at the time.

Commitments can also arise without anyone deciding. When a package, a record or an interface is relied upon by enough parties, withdrawing it has become costly for reasons that none of them chose individually, and a de facto commitment exists in a state where no commitment transition was ever licensed. This is one reason the five transitions cannot be treated as stages of a single pipeline, even though in well-run systems they often occur in the order listed.

Each of the five can occur without the next, and the next can occur without the previous. A candidate can be generated and never inscribed, inscribed and never evaluated, evaluated and never authorized, authorized and never committed to, and committed to without having been adequately evaluated or ever properly authorized. A system that treats any of these arrows as automatic has, in effect, delegated a decision to the order in which operations happen to run. The distinctions are not offered as a taxonomy for its own sake. They identify the places at which something must be added, whether evidence, a procedure, an authority or a dependency structure, for a candidate to count, and they make it possible to ask in each case which addition is missing.

It is also possible to say why transitions, and not claims, are the primitive objects of this account. A claim does not gain standing by being a claim. It gains standing when a system is permitted to move it from held to relied upon. A proof gains standing when it moves from candidate to lemma available for further use, a policy when it moves from proposal to binding, and a commit when it moves from branch to shared history. In each case the thing that changes is the permitted use of the item by others, and permitted use is a property of a move made in a particular arrangement, not of the item considered in isolation. Treating claims as the primitive hides this by suggesting that standing could be read off the content of the claim, and the third section shows by countermodel that it cannot be.

Two clarifications limit what is meant. First, private belief and unrecorded thought are not treated as transitions. A person may believe a candidate for any reason or none, and the account is silent on whether they should. The subject is what a shared record or a shared arrangement may treat as established, and for that reason the account concerns shared records and consequences and does not concern minds. Second, the five types are not claimed to be exhaustive. Repair, revocation and delegation are plausible further types, and the final section says why they are not treated here.
